One of those inputs, eligible retained income, differs from the
retained net income used in the national-bank earnings rule. Part~3 takes the
greater of two amounts: Call Report net income for the four preceding calendar
quarters, net of distributions and associated tax effects not already reflected
in that income; and the average quarterly net income over those four quarters.
Its maximum payout amount for the current quarter multiplies this eligible
income by the applicable payout ratio. The relevant capital buffers use the
previous quarter's closing ratios, so the declaration, income window and
capital-measurement date must remain distinct
\citep[12~CFR~3.11(a)(2)--(4)]{occcapitalpayout2026}.

The transactions being limited also extend beyond a cash dividend. Part~3's
distribution definition includes dividend declarations or payments on Tier~1
capital instruments and specified capital repurchases or redemptions, with
conditional exceptions for qualifying replacements
\citep[12~CFR~3.2]{occcapitaldefinitions2026}. The quarterly limit also
covers discretionary bonus payments and the creation of obligations to make
the specified payments. A bank's event history must therefore identify the
instrument, declaration, obligation and settlement, then reconcile their
regulatory treatment; a cash-dividend total cannot substitute for that history
\citep[12~CFR~3.11(a)(4)]{occcapitalpayout2026}.

A hypothetical bank in the first quarter of a year illustrates the income
distinction. Suppose each of the preceding four quarters has Call Report net
income of US\$10m. Assume its only distributions in that preceding calendar
year were US\$38m of declared dividends, with no associated tax adjustment.
Part~5 retained net income for that year is then US\$2m: US\$40m less
US\$38m. Part~3 eligible retained income is US\$10m, the greater of that
US\$2m remainder and the US\$10m quarterly average. These are two calculations
for the same hypothetical history, not interchangeable names for one balance.

Now suppose the bank falls within the amended rule's covered population and
applies that rule. Take its previous-quarter capital conservation buffer as
1.5\%, its applicable countercyclical amount as zero, its leverage buffer
standard as 1\%, and its measured leverage buffer as 0.2\%. These are supplied
example inputs, rather than estimates from a public balance sheet. The rule's Table~1
gives a 40\% payout ratio, which alone would produce US\$4m from the
US\$10m eligible income. Table~2 gives zero because the leverage buffer is
at most one quarter of its standard. The required lower ratio therefore gives
a zero table-based maximum, absent the specified OCC permission. Conversely,
a table entry saying that no payout-ratio limitation applies does not mean a
100\% ratio or permission under every other rule
\citep[12~CFR~3.11(a)(4), (c), Tables~1--2]{occcapitalpayout2026}.

